\documentclass{article}

\usepackage[margin=1in]{geometry}
\usepackage{amsmath}
\usepackage{booktabs}
\usepackage{graphicx}
\usepackage{tabularx}

\title{MNIST Validation Accuracy for Four- and Eight-Device ReRAM Synapses}
\author{}
\date{}

\begin{document}

\maketitle

\begin{table}[ht]
  \centering
  \caption{Validation accuracy at each deployment and training stage.}
  \label{tab:four-vs-eight-device-validation-accuracy}
  \begin{tabularx}{\textwidth}{@{}Xcc@{}}
    \toprule
    Stage & Four-device scheme & Eight-device ($G^+/G^-$) scheme \\
    \midrule
    Initial FP32 ReLU teacher
      & 97.26\% & 97.26\% \\
    Deployed on cohort A, before HWA
      & 61.48\% & 89.12\% \\
    After 10 epochs of cohort-A HWA/BPTT
      & 97.26\% & 97.42\% \\
    Immediately deployed on cohort B
      & 35.68\% & 42.52\% \\
    \bottomrule
  \end{tabularx}
\end{table}

\paragraph{Why accuracy drops on cohort B.}
The cohort-B values are zero-update deployment controls. The cohort-A weights
were re-encoded into common reachable windows and projected onto the measured
states of newly assigned cohort-B devices, but the individual cell states were
not subsequently recalibrated through cohort-B-specific hardware-aware
optimization before evaluation. Because each logical weight depends on
differences among several conductances, independent projection onto new device
curves changes the signed weight, cancellation error, and total passive load in
an edge- and layer-dependent way. In the four-device realization, for example,
the layerwise signed-drive RMS contracted by factors of 6.16 and 7.32. A single
scalar readout calibration cannot undo these nonuniform changes, which explains
the severe immediate loss of accuracy. The later recovery under cohort-B HWA
shows that the new devices remain capable of representing an accurate solution;
the poor deployment result primarily reflects the absence of device-specific
weight adaptation before evaluation.

\section*{Threshold-gated, decrease-only fine-tuning}

The following experiment is separate from the raw-trace HWA/BPTT comparison
above. Both schemes were initialized on cohort A using isotonic non-increasing
measured traces and a common reachable-window mapping. During fine-tuning,
each physical device could only stay at its current state or move down by one
adjacent pulse state per minibatch:
\[
  p_i \leftarrow
  \begin{cases}
    p_i+1, & \displaystyle
      \frac{\partial \mathcal{L}_{\mathrm{KL}}}{\partial G_i} > \tau_i
      \text{ and } p_i < p_{\max}, \\
    p_i, & \text{otherwise.}
  \end{cases}
\]
Here increasing the pulse index selects the next non-increasing conductance
state. The gate is strict: equality, negative gradients, and positive
gradients below the threshold all hold. The update uses neither a
learning-rate magnitude nor a digital shadow.

\begin{table}[ht]
  \centering
  \caption{Fixed raw-gradient thresholds calibrated at each scheme's shared
  isotonic initialization. Eight-device entries show $G^+/G^-$ thresholds.}
  \label{tab:four-vs-eight-gradient-thresholds}
  \small
  \resizebox{\textwidth}{!}{%
  \begin{tabular}{@{}lccc@{}}
    \toprule
    Gate & Four-device, layers 0 / 1
         & Eight-device, layer 0 $G^+/G^-$
         & Eight-device, layer 1 $G^+/G^-$ \\
    \midrule
    p90   & 494.02 / 5,849.46
          & 178.29 / 178.29
          & 12,716.26 / 12,716.63 \\
    p99   & 1,487.06 / 14,197.82
          & 559.64 / 559.64
          & 30,829.93 / 30,842.37 \\
    p99.9 & 3,078.10 / 22,323.40
          & 1,173.11 / 1,173.11
          & 52,414.81 / 52,409.63 \\
    \bottomrule
  \end{tabular}
  }
\end{table}

\begin{table}[ht]
  \centering
  \caption{Accuracy of KL-selected checkpoints after ten-epoch,
  one-pulse-down fine-tuning. Epoch 0 denotes the initialization; test
  accuracy is from a fresh held-out 10,000-example evaluation.}
  \label{tab:four-vs-eight-one-pulse-accuracy}
  \small
  \resizebox{\textwidth}{!}{%
  \begin{tabular}{@{}lccccc@{}}
    \toprule
    Gate & Four selected val. & Four test
         & Eight selected val. & Eight test
         & Selected epoch, four / eight \\
    \midrule
    Zero control & 38.32\% & 38.40\% & 48.54\% & 48.98\% & 0 / 0 \\
    p90          & \textbf{84.18\%} & \textbf{85.91\%}
                 & 82.40\% & 84.04\% & 1 / 1 \\
    p99          & 82.22\% & 83.65\%
                 & \textbf{86.08\%} & \textbf{86.72\%} & 7 / 6 \\
    p99.9        & 72.98\% & 74.38\%
                 & 81.74\% & 82.96\% & 10 / 10 \\
    \bottomrule
  \end{tabular}
  }
\end{table}

\begin{table}[ht]
  \centering
  \caption{Epoch-10 behavior. Eight-device saturation is the range across
  $G^+$ and $G^-$ in each layer.}
  \label{tab:four-vs-eight-one-pulse-saturation}
  \small
  \resizebox{\textwidth}{!}{%
  \begin{tabular}{@{}lcccc@{}}
    \toprule
    Gate & Four final val. & Eight final val.
         & Four saturation, L0 / L1
         & Eight saturation, L0 / L1 \\
    \midrule
    Zero control & 10.30\% & 8.76\%
      & 99.88\% / 99.75\% & 99.95--99.95\% / 99.80--99.85\% \\
    p90 & 13.76\% & 13.76\%
      & 93.70\% / 81.65\% & 97.11--97.13\% / 81.75--83.10\% \\
    p99 & 77.00\% & 78.34\%
      & 6.99\% / 8.80\% & 14.61--14.77\% / 15.80--17.70\% \\
    p99.9 & 72.98\% & 81.74\%
      & 4.94\% / 5.90\% & 9.10--9.24\% / 10.45--12.05\% \\
    \bottomrule
  \end{tabular}
  }
\end{table}

\paragraph{Comparison.}
The eight-device initialization is 10.22 percentage points more accurate than
the four-device initialization. The four-device p90 arm gives the stronger
p90 selected test result (85.91\% versus 84.04\%), but both p90 trajectories
are sharp early-stopping regimes and collapse when continued through epoch 10.
At p99, the eight-device scheme is stronger and stable: its selected test
accuracy is 86.72\%, compared with 83.65\% for four devices, while its final
validation accuracy is 78.34\%. The conservative p99.9 gate also favors eight
devices by 8.58 test-accuracy points. The best observed selected test result
is therefore eight-device p99 at 86.72\%, 0.81 points above the best
four-device result (p90 at 85.91\%).

All arms respected the intended physical rule: the maximum pulse-index jump
was one, no conductance increased, and no post-initialization nearest-state
projection or digital-shadow accumulation was used. These are exploratory
single-seed results. Thresholds were calibrated and checkpoints selected on
the validation split; the reported test split was held out until the selected
checkpoints were evaluated.

\end{document}
